%======================================================================
\chapter{Methodology}
\label{ch:method}

This chapter specifies what the system must do, describes the proposed method and the system at the level of modules and data flow, and defines the datasets, experimental setup, metrics, baselines, and resources used. Algorithms, equations of the models, and pseudocode are deferred to Chapter~\ref{ch:impl}; the project plan is given in Chapter~\ref{ch:planning}.

\section{Requirements Specification}
\label{sec:req-spec}

The requirements follow from the objectives in Section~\ref{sec:objectives} and from the conference reviews of Phase 1. Table~\ref{tab:fr} lists the functional requirements and Table~\ref{tab:nfr} the non-functional requirements with their quantitative targets. Every target is checked against a measured value in Table~\ref{tab:targets}.

\begin{table}[htbp]
\centering
\caption{Functional requirements}
\label{tab:fr}
\begin{tabularx}{\textwidth}{|L{1.2cm}|X|L{1.6cm}|}
\hline
\textbf{ID} & \textbf{Requirement} & \textbf{Obj.} \\ \hline
FR1 & Label the cap of every bottle from three synchronised views as one of six states and return \texttt{PASS}, \texttt{REJECT}, \texttt{UNCERTAIN}, or \texttt{INSPECTION\_FAILED}. & Primary \\ \hline
FR2 & Measure the fill ratio of a bottle from one side view and judge it against configurable underfill and overfill limits. & O2 \\ \hline
FR3 & Locate the four printed regions of the tube label, score each with its own anomaly model, and fail the label if any region fails or is missing. & O4 \\ \hline
FR4 & Store the raw frames, YOLO-format labels, verdict, and decision reason of every inspection. & O1, O5 \\ \hline
FR5 & Show live KPIs, a review queue for uncertain products, replay of stored inspections, CSV export, and run-time threshold settings in a dashboard. & O5 \\ \hline
FR6 & Switch between stations so that only one station's engines occupy GPU memory at a time. & O3, O5 \\ \hline
\end{tabularx}
\end{table}

\begin{table}[htbp]
\centering
\caption{Non-functional requirements and targets}
\label{tab:nfr}
\begin{tabularx}{\textwidth}{|L{1.2cm}|L{3.0cm}|X|}
\hline
\textbf{ID} & \textbf{Quality} & \textbf{Target} \\ \hline
NFR1 & Cap accuracy & Image-level accuracy $\geq0.93$, precision $\geq0.93$, recall $\geq0.97$ on a held-out test split. \\ \hline
NFR2 & Fail-safe decisions & False-accept rate $\leq3\%$, false-reject rate $\leq7\%$; a doubtful product is never passed automatically. \\ \hline
NFR3 & Fill accuracy & Fill-state agreement with the annotated surface $\geq0.93$. \\ \hline
NFR4 & Latency & Network time $<100$\,ms for three views; closed-loop cycle $\leq200$\,ms (300 bottles per minute). \\ \hline
NFR5 & Portability & Same code natively and in Docker; fall back to ONNX Runtime and then PyTorch when no compatible GPU is present. \\ \hline
NFR6 & Auditability & Every verdict traceable to stored evidence; engines versioned; paths set by configuration. \\ \hline
\end{tabularx}
\end{table}

\section{Proposed Methodology}
\label{sec:proposed-method}

The method replaces the broad Phase 1 design with three stations, each built around one inspection question, one set of models, and one decision rule. Three principles guided the design.
\begin{enumerate}
    \item \textbf{One station, one question.} Each station answers one inspection question (cap state, fill level, or label integrity) with its own models, decision rule, audit database, and API namespace. Stations share a server and a dashboard only for convenience, so one can be changed or removed without touching the others.
    \item \textbf{Asymmetric, fail-safe decisions.} A doubtful product is never passed automatically. Every station has four outcomes: \texttt{PASS}, \texttt{REJECT}, \texttt{UNCERTAIN} (sent to human review without ejection), and \texttt{INSPECTION\_FAILED} (the input could not be evaluated).
    \item \textbf{Evidence before verdict.} Every verdict is stored with the raw frames, the model outputs, and the reason for the decision, so any automated decision can be audited and the stored images can be re-used for training.
\end{enumerate}

\subsection{Station 1: cap inspection}
The three camera views of one bottle are resized to the network input and passed through the cap network as one batch of three. Non-maximum suppression (NMS) is applied to each view separately. The network is a YOLO11m instance-segmentation \emph{student} that was trained with feature-level distillation from a frozen YOLO11x-seg \emph{teacher} (Section~\ref{sec:kd-impl}); the teacher and the distillation projectors exist only during training. Each view is judged on its own, and the bottle verdict is \emph{defect-dominant}: a defect seen by any camera rejects the bottle. This design replaces the Phase 1 cascade of a bottle/cap detector followed by a MobileNetV3-Small crop classifier; the reasons are given in Section~\ref{sec:why-seg}.

\subsection{Station 2: fill level}
Figure~\ref{fig:fill-pipeline} shows the fill station as a multi-stage pipeline. Stage 1 localises the bottle with a YOLOv8s detector and keeps the highest-scoring bottle box. Stage 2 detects the liquid surface (meniscus) with a separate YOLOv8n detector run on the full frame, associates a surface box with the bottle, trims the bottle box to remove detector padding, and computes the fill ratio $\Rfill$ in closed form. Stage 3 maps the ratio into one of three fill classes (underfill, acceptable, overfill) and then into a gate verdict, with a review band around each limit when the meniscus detection is weak. The limits are configuration values, not learned parameters.

\begin{figure}[htbp]
\centering
\begin{adjustbox}{max width=\textwidth}
\begin{tikzpicture}[node distance=5mm, every node/.style={font=\footnotesize}]
\node[data, text width=2.0cm] (in) {Side-view\\frame $I$};
\node[stage, text width=2.4cm, right=6mm of in] (s1) {\textbf{Stage 1}\\Bottle detector\\conf $\geq 0.25$};
\node[stage, text width=2.4cm, right=6mm of s1] (s2a) {\textbf{Stage 2a}\\Meniscus detector\\conf $\geq 0.15$};
\node[stage, text width=2.4cm, right=6mm of s2a] (s2b) {\textbf{Stage 2b}\\Association +\\box trimming};
\node[stage, text width=2.4cm, below=10mm of s2b] (s2c) {\textbf{Stage 2c}\\$R_{\mathrm{fill}}$ (Eq.~\ref{eq:rfill})};
\node[dec, text width=2.6cm, left=6mm of s2c] (s3) {\textbf{Stage 3}\\3-class fill state:\\under / ok / over};
\node[dec, text width=2.4cm, left=6mm of s3] (g) {Fill gate\\(Alg.~\ref{alg:fillgate})};
\node[res, text width=2.0cm, left=6mm of g] (v) {Verdict +\\audit record};
\foreach \a/\b in {in/s1,s1/s2a,s2a/s2b,s2b/s2c,s2c/s3,s3/g,g/v} \draw[arr] (\a) -- (\b);
\draw[arr, dashed] (s1.north) -- ++(0,5mm) node[above, font=\scriptsize, align=center]{no bottle $\Rightarrow$ \texttt{INSPECTION\_FAILED}};
\end{tikzpicture}
\end{adjustbox}
\caption[Multi-stage fill-level pipeline (Station 2)]{Multi-stage fill-level pipeline (Station 2). Stage 1 localises the bottle, Stage 2 localises the meniscus and measures the ratio, and Stage 3 classifies the fill state and applies the review band.}
\label{fig:fill-pipeline}
\end{figure}

\subsection{Label prototype (O4)}
The label prototype inspects the printed label of the Himalaya Winter Defense Moisturizing Cream (50\,ml) tube. The input is a flat, unwrapped scan of about $1504{\times}8000$ pixels. Because the label wraps around a cylinder, each printed region appears twice in the scan. Figure~\ref{fig:label-pipeline} shows the three steps: (i) \emph{preprocessing and localisation}, in which the scan is converted to grayscale and four full-region templates are located by coarse-to-fine normalised cross-correlation, keeping the complete (uncut) occurrence of each region; (ii) \emph{ROI cropping}, in which each located region is cut out and rotated upright; and (iii) \emph{model routing}, in which each crop is sent to the anomaly model trained for that region. A label fails if any region scores above its calibrated threshold or if any region cannot be located.

\begin{figure}[htbp]
\centering
\begin{adjustbox}{max width=\textwidth}
\begin{tikzpicture}[node distance=5mm, every node/.style={font=\footnotesize}]
\node[data, text width=2.0cm] (scan) {Tube scan\\$1504{\times}8000$};
\node[stage, text width=2.4cm, right=6mm of scan] (pre) {Grayscale +\\coarse-to-fine\\template matching};
\node[stage, text width=2.2cm, right=6mm of pre] (crop) {ROI crop\\(complete\\occurrence) +\\rotate upright};
\node[stage, text width=2.4cm, right=10mm of crop, yshift=12mm] (r1) {ROI\_1 EfficientAD};
\node[stage, text width=2.4cm, below=1.5mm of r1] (r2) {ROI\_2 EfficientAD};
\node[stage, text width=2.4cm, below=1.5mm of r2] (r3) {ROI\_3 EfficientAD};
\node[stage, text width=2.4cm, below=1.5mm of r3] (r4) {ROI\_4 EfficientAD};
\node[dec, text width=2.2cm, right=6mm of r2, yshift=-4mm] (agg) {score $>$ $\tau_k$\\or missing\\$\Rightarrow$ FAIL};
\draw[arr] (scan) -- (pre); \draw[arr] (pre) -- (crop);
\foreach \r in {r1,r2,r3,r4} {\draw[arr] (crop.east) -- (\r.west); \draw[arr] (\r.east) -- (agg.west);}
\node[res, text width=3.0cm, right=6mm of agg] (exp) {Verdict, heatmaps, hot regions,\\JSON report (web app)};
\draw[arr] (agg) -- (exp);
\end{tikzpicture}
\end{adjustbox}
\caption[Label-defect prototype pipeline]{Label-defect prototype: template-matching localisation, ROI cropping, and routing of each crop to its own anomaly model.}
\label{fig:label-pipeline}
\end{figure}

The four regions are: ROI\_1, the logo and product name; ROI\_2, the ingredient line and paragraph; ROI\_3, the address, regulatory numbers, manufacturing and expiry dates, and barcode; and ROI\_4, the ``3 Way Care'' graphic.

\section{System Architecture}
\label{sec:system-architecture}

Figure~\ref{fig:system-architecture} shows the complete system. The input layer acquires frames: three synchronised views per bottle for the cap station, one side view for the fill station, and a full unwrapped scan for the label prototype. The inference layer runs each station's models as TensorRT engines (cap and fill) or ONNX models (label). The decision layer turns model outputs into verdicts with station-specific rules. The persistence and presentation layer writes each result to SQLite and to disk, pushes it to the dashboard over a WebSocket, and drives the reject actuator.

\begin{figure}[htbp]
\centering
\begin{adjustbox}{max width=\textwidth}
\begin{tikzpicture}[every node/.style={font=\footnotesize}]
% rows far enough apart that each group's heading clears the box above it
\def\rA{0} \def\rB{-4.0} \def\rC{-7.8}
% inputs
\node[data, text width=2.5cm] (c1) at (0,\rA) {3 synchronised views\\(front, rear, side)\\$2048{\times}2048$};
\node[data, text width=2.5cm] (c2) at (0,\rB) {1 side view\\of the bottle};
\node[data, text width=2.5cm] (c3) at (0,\rC) {Unwrapped tube scan\\$1504{\times}8000$ BMP};
% models
\node[stage, text width=3.0cm] (m1) at (4.0,\rA) {Distilled YOLO11m-seg\\TRT FP16, batch 3};
\node[stage, text width=3.0cm] (m2) at (4.0,\rB) {Bottle det.\ (YOLOv8s) +\\meniscus det.\ (YOLOv8n)\\TRT FP16, batch 1};
\node[stage, text width=3.0cm] (m3) at (4.0,\rC) {Template matching +\\4 EfficientAD models\\(ONNX)};
% decisions
\node[dec, text width=2.7cm] (d1) at (8.0,\rA) {Asymmetric multi-view\\quality gate};
\node[dec, text width=2.7cm] (d2) at (8.0,\rB) {Fill ratio $R_{\mathrm{fill}}$\\+ fill gate};
\node[dec, text width=2.7cm] (d3) at (8.0,\rC) {Per-ROI threshold,\\any-fail rule};
% outputs
\node[res, text width=2.9cm] (act) at (12.2,0.2) {Reject actuator signal};
\node[res, text width=2.9cm] (db) at (12.2,-2.0) {SQLite (WAL) audit,\\raw frames, YOLO labels};
\node[res, text width=2.9cm] (ui) at (12.2,-4.6) {React dashboard\\(REST, WebSocket)};
\node[res, text width=2.9cm] (web) at (12.2,\rC) {Upload web app:\\verdict, heatmaps, JSON};
\foreach \a/\b in {c1/m1,c2/m2,c3/m3,m1/d1,m2/d2,m3/d3,db/ui,d3/web} \draw[arr] (\a) -- (\b);
\draw[arr] (d1.east) -- (act.west);
\draw[arr] (d1.east) -- (db.west);
\draw[arr] (d2.east) -- (db.west);
\draw[arr] (d2.east) -- (act.west);
\begin{scope}[on background layer]
\node[grp, fit=(c1)(m1)(d1), label={[font=\footnotesize\bfseries, anchor=south west]north west:Station 1: cap}] {};
\node[grp, fit=(c2)(m2)(d2), label={[font=\footnotesize\bfseries, anchor=south west]north west:Station 2: fill level}] {};
\node[grp, fit=(c3)(m3)(d3), label={[font=\footnotesize\bfseries, anchor=south west]north west:Label prototype (O4)}] {};
\end{scope}
\end{tikzpicture}
\end{adjustbox}
\caption[Overall VisionFood QAI architecture]{Overall VisionFood QAI architecture. Each row is an independent station with its own models and decision rule; the cap and fill stations share the audit and dashboard layer, and the label prototype runs as a separate upload-and-inspect application. Blue: inference; orange: decision; green: outputs.}
\label{fig:system-architecture}
\end{figure}

Figure~\ref{fig:deploy} shows how the cap and fill stations are served. A FastAPI application exposes Station 1 under \texttt{/api/*} and Station 2 under \texttt{/api/fill/*}, with one WebSocket stream per station. Each station is owned by a service object that holds the engine, the gate, the audit logger, and the overlay renderer. Only one station's engines occupy GPU memory at a time: the cap engine is loaded at start-up, the fill engines are loaded on demand, and a station switch in the dashboard unloads one set before loading the other. The backend and an nginx container serving the built dashboard are started together with Docker Compose; the same backend can also run natively and serve the built dashboard itself.

\begin{figure}[htbp]
\centering
\begin{adjustbox}{max width=\textwidth}
\begin{tikzpicture}[node distance=5mm, every node/.style={font=\footnotesize}]
\node[blk, fill=gray!10, text width=2.6cm] (br) {Operator browser\\React dashboard};
\node[blk, fill=gray!10, text width=2.4cm, right=16mm of br] (ng) {nginx container\\ports 80 / 5173};
\node[stage, text width=3.4cm, right=18mm of ng] (api) {FastAPI backend container\\(CUDA 12.4, port 8000)};
\node[stage, text width=3.0cm, below left=10mm and -14mm of api] (s1) {PipelineService\\(cap; eager load)};
\node[stage, text width=3.0cm, below right=10mm and -14mm of api] (s2) {FillPipelineService\\(fill; lazy load)};
\node[data, text width=3.0cm, below=8mm of s1] (e1) {Cap TRT engine\\\texttt{best\_<ver>.engine}};
\node[data, text width=3.0cm, below=8mm of s2] (e2) {Fill TRT engines\\bottle + water surface};
\node[res, text width=3.0cm, below=8mm of e1] (d1) {\texttt{inspections.db}\\+ bottle bundles};
\node[res, text width=3.0cm, below=8mm of e2] (d2) {\texttt{fill\_level.db}\\+ bottle bundles};
\node[blk, fill=yellow!10, text width=2.6cm, left=10mm of e1] (gpu) {Single GPU:\\one station's engines\\resident at a time};
\draw[arr, <->] (br) -- node[above, font=\scriptsize]{HTTP} (ng);
\draw[arr, <->] (ng) -- node[above, font=\scriptsize]{REST, WS} (api);
\draw[arr] (api) -- (s1); \draw[arr] (api) -- (s2);
\draw[arr] (s1) -- (e1); \draw[arr] (s2) -- (e2);
\draw[arr] (e1) -- (d1); \draw[arr] (e2) -- (d2);
\draw[dashed] (gpu) -- (e1);
\end{tikzpicture}
\end{adjustbox}
\caption[Deployment architecture]{Deployment architecture. Two Docker containers (dashboard and backend) share bind-mounted model and audit folders with the host. Station switching unloads one engine set before loading the other.}
\label{fig:deploy}
\end{figure}

\section{Dataset}
\label{sec:dataset}
\label{sec:data-sources}

Table~\ref{tab:datasets} lists every dataset used; Appendix~\ref{appendix:dataset} describes their collection, composition, annotation format, limitations, and availability in detail. The cap set is public and was used as exported, with its own training, validation, and test splits. The fill-level detectors were trained on merged public sets. The label scans were supplied by an industrial partner. The code and data are kept in the following locations:
\begin{itemize}
    \item Repository: \url{https://github.com/ShreyasBairyKS/Capstone-Project} (cap and fill stations, dashboard, deployment).
    \item Label prototype branch: \url{https://github.com/ShreyasBairyKS/Capstone-Project/tree/himalaya-label-detection}.
    \item \sloppy Local project folder \texttt{E:\textbackslash P-25 Vision Food ai}, which holds the NSC scans in \texttt{dataset\textbackslash NSC} and the label-prototype working tree in \texttt{himalaya\_\allowbreak label\_\allowbreak detection}.
\end{itemize}
The NSC scans are not redistributed because the product images belong to the manufacturer.

\begin{table}[htbp]
\centering
\caption{Datasets used in Phase 2}
\label{tab:datasets}
\begin{tabularx}{\textwidth}{|L{2.9cm}|L{2.7cm}|X|L{2.4cm}|}
\hline
\textbf{Dataset} & \textbf{Used for} & \textbf{Content and split} & \textbf{Licence / source} \\ \hline
\texttt{bottle\_cap\_\allowbreak sdp} v7~\cite{roboflowcap} & Station 1 teacher and student & 6 cap classes, polygon masks; 7,458 / 712 / 354 images (train / valid / test), $2048{\times}2048$ & Roboflow Universe, CC BY 4.0 \\ \hline
Merged bottle/cap sets & Station 2 bottle detector & 4 public sources; 2,286 / 417 / 280 images; 2,133 / 774 / 356 bottle boxes & Public, CC BY 4.0 \\ \hline
\texttt{fill-level} v1~\cite{fillds1} + \texttt{bottle-\allowbreak liquid-\allowbreak level} v2~\cite{fillds2} & Station 2 meniscus detector & Single class \texttt{water\_\allowbreak surface}; 545 / 31 / 41 images; 1,683 / 54 / 76 boxes & Roboflow Universe, CC BY 4.0 \\ \hline
Four-class cap classification set & Model-selection benchmark (Table~\ref{tab:classifiers}) & 7,540 images; 308-image test split & Earlier project data \\ \hline
Himalaya NSC & Label prototype (O4) & 120 unwrapped tube scans, $\approx1504{\times}8000$ BMP: 78 good, 42 defective & Industrial partner (local copy) \\ \hline
\end{tabularx}
\end{table}

\section{Experimental Setup and Evaluation Metrics}
\label{sec:setup}

\subsection{Experimental Configuration}
\label{sec:exp-config}
Training ran on an NVIDIA RTX A5000 (24\,GB). All latency figures were measured inside the production Docker image on one workstation (RTX A5000, Intel Xeon W-2295 with 36 threads, CUDA 12.1 PyTorch build 2.4.1, TensorRT 10.16.1, Ultralytics 8.4.192, ONNX Runtime 1.31.1 with the CPU provider). Each GPU configuration was warmed up with 5 calls and timed over 100 calls; medians are reported. Station 1 accuracy is reported for both the trained FP32 checkpoint and the deployed FP16 engine. The cap models were selected on the 712-image validation split and reported on the 354-image test split, which was used neither for training nor for model selection or threshold setting. The closed-loop stress test drives the complete Station 1 path for 500 unpaced cycles and a second run of 300 cycles paced at 150 bottles per minute. Label-prototype results come from the experiments on the \texttt{himalaya-label-detection} branch and are reported with the protocol under which each number was obtained.

\subsection{Evaluation Metrics}
\label{sec:metrics}
With $TP$, $FP$, $FN$, and $TN$ the counts of true positives, false positives, false negatives, and true negatives,
\begin{equation}
    P=\frac{TP}{TP+FP},\qquad R=\frac{TP}{TP+FN},\qquad F_1=\frac{2PR}{P+R}.
    \label{eq:prf}
\end{equation}
\noindent\textbf{Detection and segmentation.} A detection is a true positive when its intersection over union (IoU) with an unmatched ground-truth instance of the same class reaches a threshold $t$. The average precision $\mathrm{AP}_t$ is the area under the precision--recall curve at IoU $t$; $\mAP@0.5$ averages it over classes at $t=0.5$, and $\mAP@0.5{:}0.95$ also averages over $t\in\{0.50,0.55,\dots,0.95\}$, which rewards tight boxes and masks. Mask AP uses mask IoU; box AP uses box IoU.

\noindent\textbf{Product-level errors.} For the cap gate a defective cap is the positive class. The false-accept rate (defect passed) and false-reject rate (good product rejected) are
\begin{equation}
    \mathrm{FAR}=\frac{\#\{\text{defective}\rightarrow\texttt{PASS}\}}{\#\text{defective}},\qquad
    \mathrm{FRR}=\frac{\#\{\text{good}\rightarrow\texttt{REJECT}\}}{\#\text{good}}.
    \label{eq:farfrr}
\end{equation}
When no error is observed in $n$ trials, the exact Clopper--Pearson~\cite{clopper1934} 95\% upper bound on the rate is $1-0.025^{1/n}$. The review rate is the share of products routed to \texttt{UNCERTAIN}.

\noindent\textbf{Agreement under imbalance.} Macro-F1 averages $F_1$ over classes; Cohen's $\kappa=(p_o-p_e)/(1-p_e)$ corrects observed agreement $p_o$ for chance agreement $p_e$~\cite{cohen1960}; and the Matthews correlation coefficient~\cite{chicco2020} uses all four cells of the confusion matrix.

\noindent\textbf{Fill level.} For an image with ratio $\hat R$ from the predicted surface and $R$ from the annotated surface, the error is $|\Delta R|=|\hat R-R|$; we report its median, mean (MAE), 90th percentile, and maximum, and the share of images whose fill state (Eq.~\eqref{eq:fillstate}) agrees.

\noindent\textbf{Latency.} Per-call network, pre- and post-processing, and wall time are reported as medians; closed-loop stage times as the 50th, 95th, and 99th percentiles (p50, p95, p99).

\subsection{Baselines for Comparison}
\label{sec:baselines}
Each result is compared with the baseline that isolates one design decision:
\begin{itemize}
    \item \textbf{Teacher YOLO11x-seg} for the distilled student: same data, same validation, FP32 (Table~\ref{tab:kd}).
    \item \textbf{FP32 PyTorch checkpoint} for the deployed FP16 TensorRT engine: same test split (Table~\ref{tab:seg}).
    \item \textbf{PyTorch FP32 on GPU and ONNX Runtime on CPU} for TensorRT latency (Table~\ref{tab:latency}).
    \item \textbf{Phase 1 detector plus MobileNetV3-Small classifier}, and five whole-image classifiers (MobileNetV3-L, EfficientNetV2-S, ShuffleNetV2, ResNet-50, DeiT-Tiny), for the choice of a segmentation detector (Tables~\ref{tab:classifiers} and \ref{tab:phase-compare}).
    \item \textbf{A MobileNetV3 fill-level classifier} for the geometric fill rule (Section~\ref{sec:fill-impl}).
    \item \textbf{Colour-trained EfficientAD models} for the grayscale label models (Table~\ref{tab:label-colour}).
\end{itemize}

\section{Software and Hardware Requirements}
\label{sec:requirements}

Tables~\ref{tab:software-req} and \ref{tab:hardware-req} list the resources used. Ultralytics was chosen because one library covers training, segmentation, and export for YOLOv8 and YOLO11; TensorRT because the target is an NVIDIA GPU; FastAPI because it serves REST and WebSocket endpoints from one asynchronous process; and SQLite because a single-file database with write-ahead logging needs no separate server on the line.

\begin{table}[htbp]
\centering
\caption{Software requirements}
\label{tab:software-req}
\begin{tabularx}{\textwidth}{|L{3.6cm}|X|L{3.9cm}|}
\hline
\textbf{Category} & \textbf{Tool / library} & \textbf{Version} \\ \hline
Language & Python & 3.12 (host); 3.11 (container) \\ \hline
Deep learning & PyTorch (CUDA 12.1 build) & 2.4.1 \\ \hline
Detection / segmentation & Ultralytics & 8.4.192 (deployment) \\ \hline
Inference engine & NVIDIA TensorRT & 10.16.1 \\ \hline
Portable runtime & ONNX, ONNX Runtime & 1.16+, 1.31.1 \\ \hline
Anomaly detection & Anomalib (EfficientAD-M) & 1.1.0 \\ \hline
Computer vision & OpenCV & 4.9--4.10 \\ \hline
Backend & FastAPI, Uvicorn, Pydantic & 0.111+, 0.30+, 2.8+ \\ \hline
Frontend & React, Vite, Recharts & 19.2, 8.3, 3.10 \\ \hline
Database & SQLite (WAL journal) & bundled with Python \\ \hline
Containers & Docker Compose, nginx, NVIDIA Container Toolkit & --- \\ \hline
Base image & \texttt{nvidia/cuda:12.4.1-\allowbreak cudnn-\allowbreak runtime-\allowbreak ubuntu22.04} & 12.4.1 \\ \hline
Annotation, version control & Roboflow, LabelImg, Git / GitHub & --- \\ \hline
\end{tabularx}
\end{table}

\begin{table}[htbp]
\centering
\caption{Hardware requirements}
\label{tab:hardware-req}
\begin{tabularx}{\textwidth}{|L{3.6cm}|X|}
\hline
\textbf{Component} & \textbf{Specification (measurement workstation)} \\ \hline
GPU & NVIDIA RTX A5000, 24\,GB VRAM \\ \hline
CPU & Intel Xeon W-2295, 18 cores / 36 threads \\ \hline
System RAM & 128\,GB (training workstation) \\ \hline
Cameras (Station 1) & Three synchronised views per bottle; $2048{\times}2048$ capture \\ \hline
Cameras (Station 2) & One side view per bottle \\ \hline
Label scanner (O4) & Line-scan or flatbed capture of the unwrapped tube, $\approx1504{\times}8000$ \\ \hline
OS & Windows 11 host with Docker (Linux containers) \\ \hline
\end{tabularx}
\end{table}
